% =============================================================================
% RÉSUMÉ — français (≈ 350 mots, 1/2 page)
% =============================================================================
\chapter*{R\'esum\'e}
\addcontentsline{toc}{chapter}{R\'esum\'e}

% TODO : finaliser après rédaction des chapitres résultats.

Les syst\`emes multi-agents bas\'es sur les grands mod\`eles de langage (LLM)
promettent d'orchestrer des workflows logiciels complexes, mais les architectures
hi\'erarchiques actuelles (planificateur--ex\'ecuteurs, superviseurs centralis\'es)
souffrent de modes d'\'echec syst\'emiques document\'es par le \emph{framework MAST}
(Cemri, 2025) et d'une \'erosion des gains avec les mod\`eles fronti\`eres
(Gao, 2025). Ce m\'emoire explore une alternative inspir\'ee de la
\textbf{stigmergie}, le principe de coordination indirecte par modification d'un
m\'edium partag\'e th\'eoris\'e par Grass\'e (1959) et Heylighen (2016) et appliqu\'e
ici aux agents LLM. La probl\'ematique pos\'ee est~: \emph{comment des agents LLM
peuvent-ils coordonner des t\^aches complexes par modification d'un m\'edium
stigmergique partag\'e plut\^ot que par communication directe ou contr\^ole
hi\'erarchique~?}

L'\'etude suit le paradigme \emph{Design Science Research} (Hevner, 2004~;
Peffers, 2007). Cinq objectifs de conception (OC1 \`a OC5) guident la
construction de l'artefact \textbf{StigmergiAgentic}, un framework Python
articulant un substrat de coordination (markers SQLite avec audit
JSONL append-only), un moteur de d\'ecision \`a pression ACO, des
m\'ecanismes d'apprentissage (\emph{lessons}, promotion en \emph{skills},
persistance \emph{cross-run} des protocoles), et une couche de gouvernance
(\emph{guardrails}, conformit\'e Article~14 EU AI Act).

L'artefact est valid\'e empiriquement sur deux benchmarks. Sur
\emph{TravelPlanner} (180 requ\^etes de planification sous contraintes), la
configuration stigmergique C3 atteint un \emph{final pass rate} de 21,1\,\%
avec Gemma et 22,2\,\% avec DeepSeek, surpassant significativement le baseline
\emph{planner-executor} (12,2\,\%) et marginalement les baselines solo. Sur
\emph{MigrationBench} (Java~8 \`a Java~17), une \emph{repair colony} avec
boucle ferm\'ee \emph{inspect/localize/propose/apply/build/classify/repair}
d\'emontre une trajectoire d'ing\'enierie DSR rigoureuse. Un panel
d'experts confirme la valeur op\'erationnelle de la trace d'audit pour la
conformit\'e r\'eglementaire.

Six principes de conception g\'en\'eralisables sont formul\'es au format
Gregor (2020), au premier rang desquels la \emph{primitive marker comme
contrat unique de coordination} et la \emph{gouvernance par le m\'edium}.
Les contributions th\'eoriques \'etendent la th\'eorie de la coordination de
Malone et Crowston (1994) aux \'ecologies agentiques contemporaines.

\vspace{0.3cm}
\noindent\textbf{Mots-cl\'es}~: stigmergie, syst\`emes multi-agents, grands
mod\`eles de langage, Design Science Research, coordination organisationnelle,
migration de code, gouvernance algorithmique, EU AI Act, MAST.

\cleardoublepage
